\subsection{Zuckerman derived functor module}

As an application of Corollary~\ref{cor:IrreducibleVermaQuasiAbelian},
we consider the branching problem of Zuckerman derived functor modules induced from
quasi-abelian parabolic subalgebras.

Let $(\lie{g}, K)$ be a pair with semisimple $\lie{g}$, and $(\lie{g'}, K')$ its subpair with reductive $\lie{g'}$ in $\lie{g}$.
Suppose that there exists an involution $\theta$ of $\lie{g}$ such that $\lie{k} = \lie{g}^\theta$ and $\lie{k'} = (\lie{g'})^\theta$,
and there exists a connected reductive algebraic group $G$ with the Lie algebra $\lie{g'}$ such that $(\lie{g}, G')$ is a pair.
Fix $H \in \lie{k}$ semisimple in $\lie{g}$ with real eigenvalues
and define a parabolic subalgebra $\lie{q}\coloneq \lie{q}(H)$.
Then $\lie{q}$ is $\theta$-stable.
Suppose that $\lie{q}$ is weakly $\lie{g'}$-compatible, and set $\lie{q'}\coloneq \lie{g'}\cap \lie{q}$.
We use the notation $\lie{u}, \lie{l}, \lie{u'}, \lie{l'}, \ldots$ as in Section~\ref{subsect:GeneralizedVerme}.

Write $L_K$ for the centralizer of $H$ in $K$ and set $L'_K\coloneq K'\cap L_K$.
Fix $\theta$-stable Cartan subalgebras $\lie{t'}$ and $\lie{t}$ of $\lie{l'}$ and $\lie{l}$, respectively, such that $\lie{t'}\subset \lie{t}$.
Let $\rhog{\lie{u}}$ denote half the sum of roots in $\Delta(\lie{u}, \lie{t})$.

In this subsection, we consider Zuckerman derived functor modules defined by
\begin{align*}
	\lmod{g}{q}{i}(V)\coloneq \Dzuck{K}{L_K}{i}\bigl(\ind{g}{q}(V)\bigr)
\end{align*}
for an $(\lie{l}, L_K)$-module $V$.
Then \smash{$\lmod{g}{q}{i}(V)$} is a $(\lie{g}, K)$-module.
In~\cite{KnVo95_cohomological_induction}, this module is defined by the Bernstein functor $\Pi$,
and our parametrization is so-called unnormalized version, which is written as $\!^{\mathrm{u}}\mathcal{L}$.

Set $S\coloneq \dim(\lie{u}\cap \lie{k})$.
For an $(\lie{l}, L_K)$-module $F$ with infinitesimal character $\lambda$, we say that $\lambda$ (or $V$) is in the good range if $\lambda$ satisfies $\Re(\lambda + \rhog{\lie{u}}, \alpha) < 0$ for any $\alpha \in \Delta(\lie{u}, \lie{t})$.
The following fact is a fundamental result about $\lmod{g}{q}{i}(V)$ (see~\cite[Theorem~0.50]{KnVo95_cohomological_induction}).

\begin{Fact}\label{fact:ZuckermanDerivedFunctorModule}
	Let $V$ be an irreducible $(\lie{l}, L_K)$-module with infinitesimal character $[\lambda]$.
	\begin{enumerate}\itemsep=0pt
		\item[$1.$] \smash{$\lmod{g}{q}{i}(V)$} has the infinitesimal character $[\lambda+\rhog{\lie{u}}]$.
		\item[$2.$] If $\lambda$ is in the good range, \smash{$\lmod{g}{q}{i}(V)$} is zero for $i\neq S$
		and non-zero irreducible for $i=S$.
	\end{enumerate}
\end{Fact}

\begin{Proposition}\label{prop:UniqueZuckermanModule}
	Let $F_i$, $i=1,2$, be a finite-dimensional irreducible $(\lie{l}, L_K)$-module with infinitesimal character $[\lambda_i]$ in the good range.
	If \smash{$\lmod{g}{q}{S}(F_1) \simeq \lmod{g}{q}{S}(F_2)$}, then $F_1 \simeq F_2$ holds.
\end{Proposition}

\begin{proof}
	Set $\lie{t}^*_\RR\coloneq \mathrm{span}_{\RR}\Delta(\lie{g}, \lie{t})$.
	Then the symmetric form $(\cdot, \cdot)$ is an inner product on $\lie{t}^*_\RR$
	and $\lie{t}^*_\RR$ is stable under the action of the Weyl group $W_G$ of $\lie{g}$.
	For $\mu \in \lie{t}^*$, we denote by $\Re(\mu)$ the real part of $\mu$ with respect to the real form $\lie{t}^*_\RR$.

	Fix a set $\Delta^+(\lie{l}, \lie{t})$ of positive roots of $\Delta(\lie{l}, \lie{t})$.
	We may assume that $\lambda_1$ and $\lambda_2$ are anti-dominant with respect to $\Delta^+(\lie{l}, \lie{t})$.
	By Fact~\ref{fact:ZuckermanDerivedFunctorModule}, $\lmod{g}{q}{S}(F_i)$, $i=1,2$, has the infinitesimal character~$[\lambda_i + \rhog{\lie{u}}]$.
	By assumption, $\Re(\lambda_i + \rhog{\lie{u}})$, $i = 1,2$, is regular and anti-dominant with respect to~${\Delta^+(\lie{l}, \lie{t}) \cup \Delta(\lie{u}, \lie{t})}$.

	Since $\lmod{g}{q}{S}(F_1) \simeq \lmod{g}{q}{S}(F_2)$, there exists $s \in W_G$ such that $s(\lambda_1 + \rhog{\lie{u}}) = \lambda_2 + \rhog{\lie{u}}$.
	By the anti-dominance, $s$ is identity and hence we obtain $\lambda_1 = \lambda_2$.
\end{proof}

In a special case, the functor \smash{$\Dzuck{M}{M \cap L_K}{i}$} can be computed by \smash{$\Dzuck{K'}{K' \cap L_K}{i}$}.
It is a generalization of~\cite[Lemma~7]{GrWa00}, which is for discrete series representations.

\begin{Lemma}\label{lem:ZuckermanKtoH}
	Let $V$ be an $(\lie{l}, L_K)$-module and $M$ a connected reductive subgroup of $K$ acting on $K/L_K$ transitively.
	Fix $i \in \NN$.
	Then there exists an isomorphism
	\begin{align*}
		\lmod{g}{q}{i}(V) \simeq \Dzuck{M}{M \cap L_K}{i}\bigl(\ind{g}{q}(V)\bigr)
	\end{align*}
	of $(\lie{g}, M)$-modules.
\end{Lemma}

\begin{proof}
	By the construction of \smash{$\Dzuck{K}{L_K}{i}$} (see before Proposition~\ref{prop:KInvZuckerman}), we have
	\begin{align*}
		\lmod{g}{q}{i}(V) \simeq H^i\bigl(\lie{k}, L_K; \ind{g}{q}(V)\otimes \rring{K}\bigr).
	\end{align*}
	Recall that $H^i\bigl(\lie{k}, L_K; \ind{g}{q}(V)\otimes \rring{K}\bigr)$ is the cohomology of the complex given by
	\begin{align*}
		C^j\bigl(\lie{k}, L_K; \ind{g}{q}(V)\otimes \rring{K}\bigr) = \Hom_{L_K}\bigl(\wedge^j (\lie{k}/\lie{l}_K),\ind{g}{q}(V)\otimes \rring{K}\bigr).
	\end{align*}
	See~\cite[Section~I.8]{BoWa00_continuous_cohomology}.
	Since $M$ acts on $K/L_K$ transitively, the restriction map induces a bijection
	\begin{align*}
		(\rring{K}\otimes W)^{L_K} \xrightarrow{\simeq} (\rring{M} \otimes W)^{M\cap L_K}
	\end{align*}
	for any $L_K$-module $W$.
	Note that $M\cap L_K$ is reductive since $M/(M\cap L_K) \simeq K/L_K$ is affine.
	Hence we have an isomorphism
	\begin{align*}
		&\Hom_{L_K}(\wedge^j \bigl(\lie{k}/\lie{l}_K),\ind{g}{q}(V)\otimes \rring{K}\bigr) \\
		&\qquad{}\simeq \Hom_{M\cap L_K}\bigl(\wedge^j \bigl(\lie{m}/\lie{m}\cap \lie{l}_K\bigr), \ind{g}{q}(V)\otimes \rring{M}\bigr)
	\end{align*}
	of vector spaces.
	It is easy to see that this isomorphism induces an isomorphism of complexes.
	Taking the cohomology, we obtain the lemma.
\end{proof}

\begin{Example}
	If $K$ is simple, there are few tuples $(K, M, L_K)$ such that $M$ acts on $K/L_K$ transitively,
	\begin{align*}
		(K, M, L_K)={}& (\lieSL(2n,\CC), \lieSp(n,\CC), \lieGL(2n-1, \CC)), \\
		&(\lieSO(2n,\CC), \lieSO(2n-1,\CC), \lieGL(n,\CC)), \\
		&\bigl(\lieSO(7,\CC), \lieG, \lieSO(5,\CC)\times \CC^\times\bigr)
	\end{align*}
	are examples.
	See~\cite[Section 3 and Example 3.7]{Ko11}.
\end{Example}

We shall consider the branching problem of \smash{$\lmod{g}{q}{S}(F)$}.
To apply Corollary~\ref{cor:IrreducibleVermaQuasiAbelian} to \smash{$\lmod{g}{q}{S}(F)$}, we need an additional assumption.
We assume that $K'$ acts on $K/L_K$ transitively and $\lie{q}$ is quasi-abelian with respect to $\lie{g}'$.

\begin{Theorem}\label{thm:DerivedFunctorModuleQuasiAbelian}
	Let $F$ be a finite-dimensional irreducible $(\lie{l}, L_K)$-module in the good range.
	\begin{enumerate}\itemsep=0pt
		\item[$1.$] \smash{$\lmod{g}{q}{S}(F)|_{(\lie{g}', K')}$} is decomposed into a direct sum of irreducible modules of the form \smash{$\lmod{g'}{q'}{S}(F')$} with finite-dimensional irreducible \smash{$\bigl(\lie{l}', L'_K\bigr)$}-module $F'$ in the good range.
		\item[$2.$] \smash{$\ind{g}{q}(F)|_{\lie{g'}, L'_K}$} is decomposed into a direct sum of irreducible modules of the form \smash{$\ind{g'}{q'}(F')$} with finite-dimensional irreducible \smash{$\bigl(\lie{l}', L'_K\bigr)$}-module $F'$ in the good range.
		\item[$3.$] For any finite-dimensional irreducible \smash{$\bigl(\lie{l'}, L'_K\bigr)$}-module $F'$ in the good range, \smash{$\Dzuck{K'}{L'_K}{S}$} induces an isomorphism
		\begin{align*}
			\Hom_{\lie{g}', L'_K}\bigl(\ind{g'}{q'}(F'), \ind{g}{q}(F)|_{\lie{g'}, L'_K}\bigr) \xrightarrow{\simeq} \Hom_{\lie{g}', K'}\bigl(\lmod{g'}{q'}{S}(F'), \lmod{g}{q}{S}(F)|_{\lie{g'}, K'}\bigr)
		\end{align*}
		of \smash{$\univ{g}^{G'}$}-modules, and
		each \smash{$\univ{g}^{G'}$}-module \smash{$\Hom_{\lie{g}', K'}\bigl(\lmod{g'}{q'}{S}(F'), \lmod{g}{q}{S}(F)|_{\lie{g'}, K'}\bigr)$} is ir\-re\-du\-cible or zero.
	\end{enumerate}
\end{Theorem}

\begin{Remark}
	The abstract branching law of \smash{$\lmod{g}{q}{S}(F)|_{(\lie{g}', K')}$} is known in~\cite[Corollaries 5.7 and 5.8]{Os24} in the Grothendieck group level.
	His result for $A_{\lie{q}}(\lambda)$ (i.e., $F$ is a character) is done for~$\lambda$ in the weakly fair range.
	His results are proved under weaker assumptions than ours.
\end{Remark}

\begin{proof}[Proof of Theorem~\ref{thm:DerivedFunctorModuleQuasiAbelian}]
	The assertion~2 has been proved in Theorem~\ref{thm:QuasiAbelianGoodRange}.
	By Lemma~\ref{lem:ZuckermanKtoH}, there exists an isomorphism \smash{$\lmod{g}{q}{S}(F) \simeq \Dzuck{K'}{L'_K}{S}\bigl(\ind{g}{q}(F)\bigr)$} of $(\lie{g}, K')$-modules.
	Then we have
	\begin{align*}
		\ind{g}{q}(F)|_{(\lie{g}', L'_K), \univ{g}^{G'}} \simeq \bigoplus_{F'} \ind{g'}{q'}(F') \otimes \Hom_{\lie{g}', L'_K}\bigl(\ind{g'}{q'}(F'), \ind{g}{q}(F)\bigr),
	\end{align*}
	and hence
	\begin{align*}
		\lmod{g}{q}{S}(F)|_{(\lie{g}', K'), \univ{g}^{G'}} \simeq \bigoplus_{F'} \lmod{g'}{q'}{S}(F') \otimes \Hom_{\lie{g}', L'_K}\bigl(\ind{g'}{q'}(F'), \ind{g}{q}(F)\bigr).
	\end{align*}
	The sum is taken over all finite-dimensional irreducible \smash{$\big(\lie{l}', L'_K\big)$}-modules in the good range.
	
	By Fact~\ref{fact:ZuckermanDerivedFunctorModule}, each \smash{$\lmod{g'}{q'}{S}(F')$} is non-zero and irreducible.
	By Proposition~\ref{prop:UniqueZuckermanModule}, \smash{$\lmod{g'}{q'}{S}(F'')$} is not~isomorphic to \smash{$\lmod{g'}{q'}{S}(F')$} for $F''\not \simeq F'$ in the good range.
	This implies
	\begin{align*}
		\Hom_{\lie{g}', L'_K}\bigl(\ind{g'}{q'}(F'), \ind{g}{q}(F)|_{\lie{g'}, L'_K}\bigr) \simeq \Hom_{\lie{g}', K'}\bigl(\lmod{g'}{q'}{S}(F'), \lmod{g}{q}{S}(F)|_{\lie{g'}, K'}\bigr).
	\end{align*}
	This isomorphism is induced by the functor $\Dzuck{K'}{L'_K}{S}$.
	The irreducibility of the $\univ{g}^{G'}$-module follows from Corollary~\ref{cor:IrreducibleVermaQuasiAbelian}.
	We have proved the theorem.
\end{proof}
